\begin{abstract}
Optimal transport（OT）理論は、フランスの数学者 Gaspard Monge（1746--1818）の言葉を借りて非形式的に次のように表現できる：建設現場に積まれた大きな砂山を、シャベルを持った作業員が移動させなければならない。作業員の目標は、その砂をすべて使って、所定の形状（例えば巨大な砂の城）の目標となる砂山を築くことである。当然、作業員は総労力を最小化したいと考える。その労力は、例えばシャベル一杯の砂を運ぶ総距離や総時間として定量化される。
%
OT に関心を持つ数学者は、この問題を2つの確率分布---同じ体積を持つ2つの異なる砂山---の比較問題として定式化する。彼らは、第1の砂山を第2の砂山に変形、\emph{輸送}、または再構成するあらゆる可能な方法を考え、砂の一粒をある場所から別の場所へ移動させるのにかかるコストという「局所的」な観点を用いて、各輸送に「大域的」なコストを対応づける。数学者が関心を持つのは、その最小コストの輸送の性質と、その効率的な計算である。
%
この最小コストは分布間の距離を定義するだけでなく、確率分布の空間上に豊かな幾何学的構造をもたらす。この構造は、分布が定義される基底となる「ground」空間の重要な幾何学的性質を受け継ぐという意味で標準的（canonical）である。例えば、基底空間がユークリッド空間である場合、補間、重心、凸性、関数の勾配といった重要な概念は、OT 幾何学を備えた分布の空間へ自然に拡張される。

OT は多くの場面でさまざまな形で（再）発見されてきたため、豊かな歴史を持つ。Monge の先駆的な研究は工学的問題に動機づけられたものであったが、1920年代の Tolstoi、1940年代の Hitchcock、Kantorovich、Koopmans がロジスティクスおよび経済学における OT の重要性を確立した。Dantzig は1949年に線形計画法の枠組みの中で OT を数値的に解き、最適化における OT の確固たる基盤を与えた。OT は1990年代に解析学者、とりわけ Brenier によって再検討され、同時にコンピュータビジョンにおいて earth mover's distance の名で知られるようになった。
%
近年、大規模な問題次元にスケーリング可能な近似ソルバーの登場により、OT の普及にさらなる革命が起きた。その結果、OT はイメージングサイエンス（色やテクスチャの処理など）、グラフィクス（形状操作）、機械学習（回帰、分類、生成モデリング）におけるさまざまな問題の解決にますます活用されるようになっている。

本論文は、数値手法に重点を置いて OT をレビューし、新しいアルゴリズムの設計を導きうる OT の理論的性質を扱う。
%
特に、OT がデータサイエンスにおいて重要性を見出すことに貢献した効率的アルゴリズムの最近の波に焦点を当てる。わずか数年の間に提案された OT の多くの一般化に重要な位置を与え、それらを統計的推論、カーネル法、情報理論に由来する関連手法と結びつける。
%
すべての図は、コンパニオンウェブサイト\footnote{\url{https://optimaltransport.github.io/}}で公開されているコードを用いて再現可能である。このウェブサイトでは、書籍プロジェクト Computational Optimal Transport をホストしており、スライドや計算リソースも掲載されている。
\end{abstract}
